\section{Streaming Translation Services}


\subsection{Proposed Architecture}

The purpose of the new Streaming Translation Service (STS) design is to enable
``near real-time delivery'' of experimental data to analysis, in the form of
both Event and Histogram NeXus files, within seconds or minutes of the
completion of a run.

\begin{figure}[htb]
\centering
\includegraphics[width=5.5in]{chunky.png}
\caption{Initial ``Chunky'' Streaming Translation Architecture, Includes
generation of ``Mini-NeXus'' files for preliminary analysis and transport to
Central computing center.}
\label{fig:chunky_sts}
\end{figure}


\begin{figure}[htb]
\centering
\includegraphics[width=5.5in]{creamy.png}
\caption{Effective Potential ``Extra Creamy'' Streaming Translation
Architecture, does Not include generation of ``Mini-NeXus'' files, only produces
a Full NeXus file for ``soon after post-mortem'' analysis, via transport (or
direct writing) to the Central computing center.}
\label{fig:creamy_sts}
\end{figure}


The basic proposed architecture for this STS design
is illustrated in Figures~\ref{fig:chunky_sts} and \ref{fig:creamy_sts},
for two incremental configurations based on use of intermediate
``Mini-NeXus'' file generation. Note that the actual boundary/separation between
the Data Acquisition System (DAS) and Central Computing Center is transient
here, and is expected to move as changes to the underlying network
infrastructure solidify. Yet moving or removing that boundary does not
fundamentally change the design elements or associated development tasks; e.g.
Parts 1 \& 2 of the Streaming Translation software can be combined into a
single Translation entity on the Central Computing side, without loss of
generality.


Much existing technology can be applied to assemble this new Streaming
Translation architecture, pulling code and expertise from each of the Live
Data Processing (LDP) ``Event Catcher'', the HistoTools Event- and
Histogram-based NeXus file generation applications \& library, and the
existing static Translation Service (TS) Client and TS Monitor GUI.  


In short, the intent is to attach an LDP ``front end'' onto a HistoTools
library ``back end'' to parse the Live Data Stream coming from the DAS and
generate Live Event (and Histogram) NeXus files while the experiment is
running. Based on existing technological capabilities, this Streaming
Translation solution should be able to adequately ``keep up'' with even the
highest bandwidth SNS beamline, and effectively be ``done'' creating a given
NeXus file (whether as a single monolithic file, or via the assembly of a
sequence of Mini-NeXus files) within moments after the last data packet
arrives, signaling the end of the experiment. Then the traditional Translation
Service Client can perform its usual Archiving and Cataloging services, along
with new extensions to improve performance, such as providing multi-threaded
concurrency in the translation processing to effectively ``overlap'' independent
operations and capitalize on powerful computational resources at SNS.

\subsubsection{Implementation Plan}

The STS implementation relies heavily on completion of the
Hybrid Streaming Network Protocol,
as described in Section~\ref{network_protocol}.
While some STS architectural re-organization and development can proceed
in unison with the design of the network protocol,
some practical results will need to be available early
in the STS development.
Specifically, a ``Protocol Generator'' must be created to take existing
historical SNS data (likely in the ``Pre-NeXus'' binary format, and the
associated XML meta-data files) and produce a network stream using
the new hybrid network protocol,
for proper design and testing of STS subsystems.
All specific packet types should be well defined,
including accelerator pulses, neutron events, beam monitor events,
chopper/special events, analog/pump probe sampling events,
polarization state transitions, slow controls meta-data/motor positions
(geometry) and accelerator information.
Beyond this preliminary stand-alone network protocol generator,
eventually a full ``side-by-side'' testbed should be provided,
in a real beamline such as HYSPEC, for testing in a production
experiment environment.

The front-end of the new STS system will be derived from the
existing SNS production ``Live Data Processing'' (LDP)
Event Catcher software,
which is utilized to snoop packets from the legacy DAS live event stream
and produce transient live histogram-based NeXus files
while an experiment is running.
An optional intermediate development step for STS would be to enable a
preliminary extension of the existing LDP stand-alone subsystem to produce
live {\em Event-based} ``Mini-NeXus'' files,
in preparation for subsequent potential streaming translation service needs.
These Mini-NeXus files {\em might} prove ``interesting''
for early Live Analysis developments.

Otherwise, the eventual usage for the current LDP Event Catcher
would be as the front-end linkage between the new hybrid network protocol
and a back-end streaming translation client.
Based on the new network protocol,
a modified communication/control linkage will be incorporated
into the existing LDP network handler,
and then data can be captured and delivered appropriately
to the back-end client for processing and storing the data into NeXus files.
This will require modified handling for beamline geometry information,
which is currently extracted from a Live Geometry MySQL Database
in the existing Event Catcher software.
A migration must occur from this database infrastructure
to instead extract slow controls process variable updates
from the new network protocol packet types,
and including all experiment-related meta-data on
users, proposal ids, motor positions, etc.

This back-end streaming translation client
will itself consist of several subsystems,
also to be constructed from other existing production SNS software projects,
namely the ``HistoTools'' NeXus file generation effort
and the legacy Translation Service Client.
The HistoTools application code and library can be applied,
by a straightforward manipulation of the data input method,
to process a network stream instead of file based data.
This is especially true given the recent optimizations and
algorithmic reworking in the impending HistoTools 4.0.0 production release,
which processes input data from a ``pulse-driven'' mode
instead of an ``event-driven'' one.
While more efficient in terms of file-based processing,
this high-level data handling approach is also directly amenable
to the processing of network packets in the new data protocol.
A first step in this development will be to effectively replace
the existing LDP back-end NeXus file generation code
with linkage to the new HistoTools application/library.

An optional companion effort with the HistoTools software
would be to provide basic technology for
splitting and merging ``partial'' Mini-NeXus files,
as would be needed to cooperate with any Mini-NeXus files
generated in conjunction with the LDP Event Mini-NeXus file option,
or as part of the ``Chunky'' STS design.
As such, a file merger application/library should be developed
for incrementally concatenating Mini-NeXus Files into
single monolithic NeXus files.
This work would include a full performance evaluation
and code ``hot-spot'' optimization exploration
to determine the feasibility of this approach in the STS overall design.
For example, an optimal ``chunk size'' for Mini-NeXus files
could be determined,
based on the overheads involved in concatenating them, etc.

The second key area of functionality in the STS back-end
comes from the existing production Translation Service Client software.
This software currently manages the
notification of experiment completion,
coordinates the execution of the HistoTools software
to produce Histogram and Event-based NeXus files,
and then archives and catalogs these files in ICAT and VFS,
along with the original raw experiment data.
In the new STS design, many of these same tasks will still be required,
with the exception of the experiment completion notification,
which will now be handled via direct communication in the
hybrid network protocol,
as part of the modified STS LDP/HistoTools linkage.
Yet the Translation Service Client will still need to manage
the archiving and cataloging of experimental data,
including the captured network protocol stream,
which must also be written to local storage during processing.

Additional features must be explored for this
new STS back-end Translation service,
including incorporation of ``Event Mode'' Beam Monitor data,
and other ``Fast'' Meta-Data,
such as from Magnet Pulses on SEQUOIA
and High-Voltage/Analog Pump Probe on VULCAN.
Usage requirements must be collected to ensure proper handling
of these specialized data products.

Because the existing production Translation Service Client
is used for a variety of experimental facilities beyond SNS,
including HFIR at ORNL,
it will be necessary to generalize/separate the workflow handling
of data for SNS versus HFIR and other facilities.
A new architecture for the Translation Service Client
will be designed,
to be continued... :-D

XXX START HERE XXX
March 9, 2012: Design and Implement Streaming Translation “Part 1” Front-End Control Infrastructure, Including Workflow, Communication and Oversight with LDP and HistoTools Data Collection and Processing
March 9, 2012: Design and Implement (Borrowing Significantly from Existing Translation Client Code ☺) Streaming Translation “Part 2” Back-End Archiving, Permissions Setting and Cataloging Infrastructure, Including Integration with ICAT2/ICAT3 and Orbiter VFS, and Accounting for Computing/Network Hardware Design Choices
May 31, 2012: Extend Translation Monitor to Display End-to-End System Status and Error Logging for Entire System from DAS through Live Streaming to Analysis
May 31, 2012: Optimize Streaming Translation Internal Workflow, by Incorporating Concurrency via Multi-Threading, Using “Pool of Tasks” Algorithm Approach at Each Stage to Independently Process the Sequence of Translation Phases
May 31, 2012: Explore \& Design Auto-Recovery Approach for Coordinating Recovery from Streaming Exceptions During an Experiment
XXX START HERE XXX

Integral to the development of the Streaming Translation Service
is proper capturing, handling and reporting of
System Status and Error conditions.
To emphasize the importance of this information
in the proper administration and maintenance
of the experimental facility,
consideration of Status/Error logging will be incorporated
early in the design of the new STS system.
An architecture will be designed for culminating status/error information
from across the DAS/Translation processing pipeline,
including both hardware and software errors,
as well as inherent data-related omissions or inconsistencies.
This status/error information must also be propagated
to a modified Translation Service Monitor
for easy notification and dissemination to support staff, end-users,
instrument scientists, and instrument hall coordinators.

Additional development may be explored to enable some rudimentary level of
automated recovery for handling streaming exceptions during an experiment,
e.g. regenerating key meta-data lost in the live stream,
using stream retry requests to the Data Aggregator
(which would pull a cached version of the stream from disk for replay).

